\section{Polymer ethenolysis: less fresh catalyst during reuse}
\label{sec:polymer}

\textbf{Summary.} Conk et al.\ (2024) converted polyethylene (PE) to olefins including propylene by isomerizing ethenolysis (repeated double-bond migration and metathesis with ethylene) over WO\textsubscript{3}/SiO\textsubscript{2} (W/silica) mixed with Na/$\gamma$-Al\textsubscript{2}O\textsubscript{3} (Na/alumina). When the retained mixture is reused over successive PE charges, output falls; one 0.4\,g fresh Na/alumina addition before charge 2 restores it over three charges.\cite[Fig.~4B]{conk2024} That rescue is our benchmark; pressure effects, initiation and regeneration, and chain-population models are also prior work.\cite{conk2025thesis,hartwig2026patent,guironnet2020,chen2024,wang2022thesis} The open question is whether lower ethylene partial pressure can replace part of the fresh Na/alumina, predictably. We propose to reproduce the published reuse, confirm the reuse deficit in a semibatch reactor, which pressure control requires, cross a fresh Na/alumina dose bracket with one lower ethylene pressure, predict and test a new intermediate dose, and extend the selected policy over further charges. Success would give an operating rule that needs less fresh catalyst per unit of polymer-derived carbon at specified product quality and complete cycle time, reducing catalyst preparation, solids handling and waste while preserving polymer-carbon recovery; three successful batches would not establish economic superiority. At equal output, a 0.2\,g midpoint addition would use 25\% less fresh Na/alumina and 16.7\% less total fresh catalyst solid than the published 0.4\,g addition (arithmetic, not a measurement). Selective loss of an initiating function is one possible explanation of the reuse loss, not an established result or a prerequisite for the operating rule.

\subsection{Prior work}

Conk and colleagues converted PE and polypropylene directly to olefins including propylene; W/silica supplies olefin metathesis and Na/alumina olefin isomerization, but these roles do not show how saturated PE first becomes reactive or which function changes during reuse. The standard experiment uses 1\,g PE and 0.4\,g of each solid at 320\,$^\circ$C, with high fresh-batch output and demonstrated reuse.\cite{conk2024} In the reuse benchmark the mixture is retained whole, and charge 3 receives PE alone.\cite[Fig.~4B]{conk2024} Rescue does not show that W function is unchanged, since added Na/alumina can also alter unsaturation supply, isomerization, activation, contact, and the balance among functions.

Conk's public dissertation abstract proposes initiation by transfer dehydrogenation with ethylene as hydrogen acceptor, rather than the paper's cracking interpretation, and describes modified catalysts and dimethyl-ether-based regeneration; the body is embargoed, so no protocol or performance comparison can be verified.\cite{conk2025thesis} Patent WO2026011187A2 discloses oxidative then dimethyl ether treatment, but its Figure~21 caption names a Na--Fe formulation and the adjoining procedures are not clearly consistent about catalyst identity, so it cannot be treated as a demonstrated regeneration benchmark for every unmodified Na/W mixture.\cite{hartwig2026patent}

Guironnet and Peters' chain-population model links ethylene and reactive-chain concentrations; they describe a numerical extension to changing ethylene but display only fixed-ethylene solutions.\cite{guironnet2020} Chen, Sadow, and Peters' final 2024 model resolves finite isomerization kinetics and predicts saturation and possible ethylene inhibition, but for a molecular catalyst with initially unsaturated chains, without deactivation or ethylene gradients, and physically restricted to relatively slow isomerization (its quasi-equilibrium limit is only a numerical check). It does not predict that lowering ethylene improves aged Na/W.\cite{chen2024} Wang experimentally examined pressure reduction and catalyst supply in an earlier tandem system with different materials; its pressure comparisons also changed other conditions.\cite[ch.~4]{wang2022thesis}

\subsection{Carbon origin and catalyst demand}

Gross propylene is the wrong measure because ethylene also supplies product carbon. The source's labeled-polymer experiment shows polymer participation under its fresh-batch conditions, not a general carbon allocation for every reused inventory and pressure.\cite[Fig.~4A]{conk2024} The accepted output $Q_{\mathrm{PE}}$ is the moles of carbon from supplied PE recovered in the specified accepted product over a whole sequence. For a propylene objective, count polymer-origin carbon in propylene and report higher olefins separately. Define
\begin{equation}
 D_{\mathrm{Na}}=\frac{M_{\mathrm{Na/alumina,fresh}}}{Q_{\mathrm{PE}}},\qquad
 D_{\mathrm{W}}=\frac{M_{\mathrm{W/silica,fresh}}}{Q_{\mathrm{PE}}},\qquad
 J_{\mathrm{PE}}=\frac{Q_{\mathrm{PE}}}{T_{\mathrm{complete}}}.
 \label{eq:polymer-demand}
\end{equation}
These are catalyst-solid demands, not elemental Na or W demands; report elemental Na and W inventories and retained and discarded solids separately. Complete time includes heating, reaction, cooling, gas handling, unloading or recharge, and any regeneration.

To attribute carbon, run representative sequences with ethylene at the \emph{same enrichment throughout conditioning and every charge}, plus unlabeled reference sequences; this avoids assuming that a small labeled-PE fraction represents all chain populations. Replace the source's methane internal standard with He, measured by calibrated GC with a thermal-conductivity detector and Ar carrier (the source's flame-ionization detector cannot detect He), keeping separately calibrated hydrocarbon detection and synchronized sampling. First validate He recovery against known gas flows, separation from hydrogen, response in the actual gas matrix, and analytical access, and recheck reference behavior after the tracer change. Quantify initial catalyst carbon and other carbon inputs; appreciable unclassified carbon invalidates two-source attribution.

Where PE and ethylene are the only contributing carbon reservoirs and isotope discrimination is bounded, the polymer-origin fraction in product $j$ is
\begin{equation}
 f_{\mathrm{PE},j}=\frac{x_E-x_j}{x_E-x_{\mathrm{PE}}},\qquad
 Q_{\mathrm{PE}}=\sum_j N_{C,j}\,n_j\,f_{\mathrm{PE},j},
 \label{eq:polymer-isotope}
\end{equation}
where $x$ denotes \emph{carbon atom fraction} of $^{13}$C, $N_{C,j}$ the number of carbon atoms per product molecule, and the sum covers accepted products collected over the sequence. Derive $x_j$ from calibrated isotopologue distributions with natural-abundance and fragmentation corrections; the published singly labeled propylene abundance is not a carbon atom fraction. First validate isotopic response and bound fractionation or unequal enrichment within a source; do not force fractions into the interval zero to one.

Carbon collection covers gas, condensates, soluble and insoluble residue, and catalyst-bound carbon. Constant enrichment assigns retained carbon to its source but not to a particular PE charge, so decisions use cumulative sequence output; charge-resolved kinetics must state carryover. Credit collected products once, not also as disappearance of intermediate residue. A polymer-free blank bounds background reactions but does not reproduce a polymer-conditioned catalyst.

\subsection{Experiment 1: reproduce reuse and bracket the fresh Na/alumina dose}

With one characterized lot of clean HDPE, reproduce the catalyst preparation, one fresh time course, and the three-charge unsupplemented and once-supplemented sequences. The sealed procedure charges ethylene and methane before heating, so its nominal 15\,bar ethylene is not a controlled hot partial pressure; record fill conditions, pressure convention, and actual hot pressure and temperature history. The separate-component figure lists 10\,bar where the text and supporting procedure specify 15\,bar; follow the detailed method and record the discrepancy.\cite{conk2024,conk2024si}

Table~\ref{tab:polymer-makeup} gives the dose bracket; 0.2\,g is our chosen midpoint. Each reused arm receives its fresh Na/alumina makeup before charge 2 only and keeps the whole mixture without washing or separation. The all-fresh arm, with fresh catalyst for each of three charges, is the output reference.

\begin{table}[ht]
\centering
\caption{Fresh catalyst solids consumed over three 1\,g PE charges. The lower dose is proposed; the full once-only dose is the published comparator.}
\label{tab:polymer-makeup}
\begin{tabular}{@{}lrrr@{}}
\toprule
Policy & Na/alumina (g) & W/silica (g) & Total (g)\\
\midrule
Fresh mixture for every charge & 1.2 & 1.2 & 2.4\\
Retained, no makeup & 0.4 & 0.4 & 0.8\\
Retained, 0.2\,g before charge 2 & 0.6 & 0.4 & 1.0\\
Retained, 0.4\,g before charge 2 & 0.8 & 0.4 & 1.2\\
\bottomrule
\end{tabular}
\end{table}

With $f=Q_{\mathrm{PE,reuse}}/Q_{\mathrm{PE,allfresh}}$, total-solid productivity relative to all-fresh operation is $2f$ for the published policy and $2.4f$ for the midpoint, which gives the savings quoted in the Summary at equal output. Three doses bracket demand; they do not locate a precise optimum.

The reported restart cools and dismantles the reactor under inert conditions. A polymer-free thermal and handling cycle applied to fresh catalyst before its first PE charge tests whether interruption alone changes function (not the spent polymer environment); record solid recovery, retained carbon, exposure time, and transfer losses. Closed inert recharging, if available, is compared only after validating delivered PE mass, temperature history, and mixing; it may change more than exposure.

An added-support control (same dried support at the full supplement's solid mass) tests whether extra solids or contact explain part of the rescue; it cannot reproduce Na/alumina surface chemistry or rheology. Where possible, check agitation and particle-size sensitivity; a stirring plateau does not exclude pore-access limits. If handling dominates the loss, the outcome is a handling protocol, not evidence of intrinsic, turnover-driven deactivation.

Fix in advance the required cumulative accepted output, allowed product distribution, and complete-time ceiling, and report each charge's output so early production cannot hide a stalled final charge.

\subsection{Experiment 2: one lower ethylene pressure crossed with dose}

The pressure intervention needs a semibatch reactor with continuous gas flow and product removal, whereas the benchmark is sealed batch. First establish the fresh--retained difference and the direction of Na rescue in semibatch mode; without an important reuse deficit there, stop testing pressure as a rescue.

Lower ethylene partial pressure by inert replacement at fixed total pressure, inlet molar flow, temperature, PE charge, catalyst inventory, and agitation. As a starting contrast, halve the reference inlet ethylene fraction, keep a constant carbon-free tracer fraction, and make up the balance with Ar, within the validated gas-delivery range. Measure outlet composition and actual molar flow rather than inferring reactor ethylene partial pressure from the cylinder blend, and record absolute versus gauge pressure. The source semibatch procedure's 19\,bar gauge total pressure is not interchangeable with its sealed-batch ethylene fill pressure.\cite{conk2024si}

Match all history through charge 1 at the reference condition, then randomize retained inventories across both pressures and the Experiment 1 doses (0, 0.2, 0.4\,g Na/alumina, added once before charge 2): six conditions. Test the nonzero doses whatever the zero-dose result, because pressure and makeup can interact. Fresh mixtures at both pressures show whether a response is specific to retained catalyst. The contrast is an operating intervention on a changing mixture, not an intrinsic reaction order. The source raised semibatch output while increasing pressure and flow together; that cannot establish a pressure order, but it is substantial evidence against assuming that less ethylene must help.\cite{conk2024}

Compare complete charges, not rapid pressure waveforms: the published apparatus samples every 12 minutes, and gas holdup and dissolved-product release may obscure shorter responses.\cite{conk2024si} A reference--lower--reference transient helps only if measured response times are short compared with chain evolution and catalyst change; a propylene spike from stored product is not restored production.

A useful reduced-dose policy must meet polymer-carbon attribution, full product collection, product specification, and the complete-time ceiling. Compare each reduced dose at both pressures with the full-dose, reference-pressure policy, reporting the pressure contrast within each dose. A null result at zero dose rules out only pressure rescue without makeup; a broader saving claim also needs the reduced nonzero dose. If the fresh comparison supports it, test lower ethylene from the start of all three charges as a separate comparison (its charge-1 history differs from the six-condition cross); a fixed condition that performs as well as a history-dependent switch is preferred. Count inert-gas supply, separation, and circulation; lower partial pressure at unchanged total pressure does not imply lower compression demand or stoichiometric ethylene use.

The opposite outcome is equally plausible: ethylene takes part in metathesis, may accept hydrogen during initiation, and changes swelling and gas--polymer transport, so lowering it can slow one function more than it relieves another, shift selectivity, or reduce contact. Drop the partial-substitution claim if the crossed comparisons exclude an important saving that meets the output, specification, and time requirements. Either way, the result does not identify which function decayed, and a pressure library is not justified before an apparent improvement appears.

\subsection{Chain-state evidence and selective-initiation claims}

Product trajectories are tracked with full molecular-weight distributions, recovered fractions, and, where measurable, unsaturation. Endpoint olefin content reflects formation and consumption together, so low unsaturation does not prove slow initiation, nor does a bulk C=C fraction show how many chains carry reactive sites; heating alone changes molecular weight in the source's component experiments, so mean molecular weight alone cannot assign mechanism.\cite{conk2024} Because depressurization and two to three hours of cooling allowed isomerization and fragment recombination in Wang's tandem system,\cite[ch.~4, pp.~72--74]{wang2022thesis} compare the fastest validated product workup with a deliberately delayed workup at the same conversion window. Solvent-soluble material cannot represent carbon retained on inaccessible solids.

If pressure gives a real saving, a matched unsaturated PE and its hydrogenated counterpart can test whether feed state changes the benefit, with chain-length and branching distributions, workup residues, and metal contamination controlled and a prediction from the measured Na/W response, not the optimum of Chen et al.'s molecular Ru-catalyst model. It does not isolate initiation, because unsaturation also changes chain distribution and contact. A component-specific claim needs probes shown to be selective on the separate catalysts and supports, since small olefins may reach sites that polymer cannot or activate the material; otherwise, keep an empirical relation among history, pressure, and fresh-material demand.

\subsection{Experiment 3: predict a new dose, then test whether the saving persists}

Build the simplest predictor of accepted output at both pressures from the crossed brackets and complete time courses, keeping any pressure--dose interaction: start with interpolation over the measured makeup range and its uncertainty; do not fit separate unobserved initiation and metathesis capacities to one product trace. If output is not regular enough in dose to interpolate, choose among tested policies.

\emph{Validation.} Choose a new intermediate dose (not 0, 0.2 or 0.4\,g) only if the predictor says it meets the output and time criteria with a resolvable saving. Test it at \emph{both} pressures on three-charge sequences with matched charge-1 histories, added before charge 2 only, against the best tested fixed-pressure policy and the published full-supplement schedule in the same reactor mode. It succeeds if the lower confidence bound on accepted output meets the requirement, specification and time are met, and fresh-material demand is lower after attribution uncertainty. Pressure enables the saving only if the lower-pressure policy passes and the same dose at reference pressure fails an acceptance criterion beyond its uncertainty. If both pass, report an adequate lower dose and any separately resolved pressure benefit, without claiming pressure was needed; an unresolved reference result leaves this open. No minimum over unmeasured doses is identified.

\emph{Extension.} To test whether the saving only postpones replacement, run the selected policy and its comparator through further 1\,g PE charges toward a common cumulative accepted-output target, within maximum complete time and charge count, all fixed in advance with a per-charge replacement trigger (accepted output, quality, charge-time ceiling). At the trigger, remove the whole catalyst--residue mixture under the validated handling protocol, quantifying its solids and unprocessed carbon; add 0.4\,g fresh Na/alumina and 0.4\,g fresh W/silica; and restart the assigned pressure schedule, with once-only makeup before the new mixture's second charge and the common reference first charge where validation used it. No selective W replacement or unvalidated regeneration. Count failed charges and all removal, recharge, and reheating time; residues stay in the carbon and waste accounts and are not recovered product. Credit the common target as output and report excess separately; a policy that misses the target gets no extrapolated credit. The conclusion covers only this extension and replacement rule; a shorter service interval pays its measured replacement cost.

The scientific payoff grows if a chain or catalyst-state descriptor predicts a new history better than charge count alone and so changes a makeup decision; otherwise a reproducible fixed-pressure saving is still a useful operating result. If no validated intervention lowers demand at accepted output and time, drop the improvement claim. Neither a negative result nor fresh Na/alumina rescue identifies a unique molecular cause of reuse loss.
